% Generated from validated Article Draft Result. Do not hand-edit; revise the structured draft instead.
\section{フロンティアモデルは、会話相手からマルチモーダルな対話基盤へ}
\label{pkg:frontier-interaction}
\noindent\textbf{2024年上期のフロンティア競争では、長い文脈、画像理解、音声入出力、応答遅延が別々の機能ではなく、ひとつの対話体験を構成する層として結びつき始めた。GPT-4o、Gemini 1.5、Claude 3系を軸に、その変化と公開範囲の差を読む。}\autocite{src-28ee6f0f00f4,src-2d3a533e91db,src-e91e6f80eaac}

2024年上期を境に、フロンティアモデルの評価軸は「テキストでどれだけ正しく答えるか」だけでは捉えにくくなった。OpenAIのGPT-4o、GoogleのGemini 1.5系、AnthropicのClaude 3系はいずれも、視覚や長文脈を含む入力、より自然な応答、複数モダリティの扱いを製品・モデル系列の中心へ押し出した。重要なのは個別ベンチマークの順位より、ユーザーが触れるinteraction layer自体がモデル設計と不可分になったことである。\autocite{src-28ee6f0f00f4,src-2d3a533e91db,src-e91e6f80eaac}


ただし「multimodal」という語は同じ意味ではない。GPT-4oは統合された入出力モデルとして提示され、Gemini 1.5やClaude 3系は長文脈・画像入力を含む能力をそれぞれの提供形態で拡張した。一方、Chameleonのような研究はテキストと画像をより統一的な生成表現で扱う方向を示す。したがって、APIやUIで複数モダリティを使えることと、モデル内部でearly/native fusionを行うことは分けて読む必要がある。\autocite{src-28ee6f0f00f4,src-2d3a533e91db,src-e1259b3c0804,src-e91e6f80eaac}


長文脈も同様に、発表された最大context windowと、その時点で一般利用できた範囲は同義ではない。Gemini 1.5は2月のpreviewから5月のFlashや開発者向け更新へ進み、Claude 3系も3月のfamily launchから6月のClaude 3.5 Sonnetへ更新された。Quiet-STaRやInfini-attentionの研究は、長い文脈や内部推論をどう効率よく扱うかという技術側の探索が製品競争と並行していたことを示す。\autocite{src-12fa04e64d24,src-28ee6f0f00f4,src-2eb2353f0533,src-e91e6f80eaac}


\begin{claimboundary}
各社が示した速度、精度、長文脈性能、マルチモーダル能力の比較値は、それぞれの評価条件に基づくvendor claimである。本稿では発表・提供形態と技術的な方向性を比較するが、異なるベンチマークや製品条件を横断して独立した順位へ換算しない。\autocite{src-06ed597872be,src-28ee6f0f00f4,src-2d3a533e91db,src-a2cf0d0271ac,src-e91e6f80eaac}
\end{claimboundary}

この半年の転換は、モデルが「チャット欄の中で文章を返すもの」から、視覚・音声・大量のcontextを取り込みながら低遅延で応答する対話基盤へ近づいたことにある。2024-H1のフロンティア競争は、parameter scaleだけでなく、どの情報をどの時間感覚で扱えるかというsystem-levelの設計競争へ広がった。\autocite{src-28ee6f0f00f4,src-2d3a533e91db,src-e1259b3c0804,src-e91e6f80eaac}
